% 2-1-reviews.tex
%  Budget: 2pp.  Gate-1 map APPROVED (with swaps): jedrzejewski added as anchor,
%  mohammadi rescoped to ML-in-markets (not EPF), opening span = 25 years (sourced).
%  Cite ONLY the LaTeX .bib keys (e.g. trebbien_explainable_2023, not _nodate).
%  Synthesis, not a list of summaries; forward-references 2-2..2-9.

\section{مرورهای نظام‌مند و مطالعات کتاب‌سنجی}

پیش‌بینیِ قیمت برق در حدود \lr{25} سال گذشته به یکی از حوزه‌های فعالِ پژوهشی در
اقتصادِ انرژی بدل شده و بخشِ عمده‌ی این پژوهش‌ها \lr{—} تا نزدیک به \lr{90} درصد
\lr{—} بر بازار روز-پیش متمرکز بوده‌اند؛ بازاری که به‌عنوان نماینده‌ی متعارفِ
«قیمت برق» شناخته می‌شود \cite{jedrzejewski_electricity_2022}. این حوزه نقشی
بنیادی در تصمیم‌گیریِ بازیگرانِ بازار در محیط‌های آزادسازی‌شده و رقابتی دارد
\cite{jiang_review_2018, cerjan_literature_2013, oconnor_review_2025}، و با
گسترشِ منابعِ تجدیدپذیر و نوسان‌پذیریِ ناشی از آن، بر اهمیت و دشواریِ آن افزوده
شده است \cite{oconnor_review_2025, cerjan_literature_2013}.

مرورهای این حوزه، روش‌های پیش‌بینی را در چند خانواده‌ی اصلی دسته‌بندی می‌کنند.
یکی از نخستین مرورهای جامع، این روش‌ها را به سه گروهِ سری‌های زمانیِ تصادفی،
مدل‌های علّی، و مدل‌های مبتنی بر هوش مصنوعی تقسیم کرده و آن‌ها را بر پایه‌ی افقِ
پیش‌بینی، متغیرهای ورودی و عواملِ مؤثر بر قیمت مقایسه می‌کند
\cite{aggarwal_electricity_2009}. مرورهای متأخرتر این دسته‌بندی را گسترش داده و
روش‌های یادگیری ماشین، یادگیری عمیق، و معماری‌های ترکیبی و گروهی را نیز در بر
می‌گیرند \cite{oconnor_review_2025}. روندِ کلیِ ادبیات، حرکتی آشکار از روش‌های
آماری به سوی روش‌های مبتنی بر هوش مصنوعی را نشان می‌دهد \cite{jiang_review_2018,
oconnor_review_2025}؛ گرایشی که تنها به پیش‌بینیِ قیمت محدود نیست، چنان‌که مروری
گسترده‌تر بر کاربردهای یادگیری ماشین در مطالعاتِ بازار برق \lr{—} فراتر از
پیش‌بینیِ قیمت و شاملِ حوزه‌هایی چون طراحیِ بازار \lr{—} نیز همین گسترش را بازتاب
می‌دهد \cite{mohammadi_review_2020}. با این‌همه، مرورها تأکید می‌کنند که هیچ روشی
به‌طور مطلق برتر نیست و هر خانواده مزایا و محدودیت‌های خود را دارد
\cite{jiang_review_2018}.

در سطحِ یافته‌ها، یک مرورِ نظام‌مندِ اخیر که \lr{62} اثرِ منتشرشده میان سال‌های
\lr{2012} و \lr{2022} را بررسی کرده است، بر چند نکته‌ی مشترک انگشت می‌گذارد
\cite{chai_forecasting_2024}. نخست، لحاظ‌کردنِ عواملِ تعیین‌کننده‌ی قیمت می‌تواند
دقتِ مدل‌ها را به‌طور چشمگیری بهبود دهد، هرچند بر سرِ این‌که کدام عوامل و به چه
شماری باید در نظر گرفته شوند اجماعی وجود ندارد. دوم، بخشِ عمده‌ی پژوهش‌ها بر
پیش‌بینیِ نقطه‌ای متمرکزند، در حالی که پیش‌بینیِ بازه‌ای و احتمالاتی \lr{—} که به
دامنه و عدم‌قطعیتِ تغییراتِ قیمت پاسخ‌گوتر است \lr{—} کمتر بررسی شده است. افزون
بر این، عملکردِ نسبیِ مدل‌ها به بخشِ بازار وابسته است: مروری که بازارهای روز-پیش،
درون‌روز و تعادلی را کنارِ هم می‌سنجد، برتریِ مدل‌های ترکیبی و گروهی را در بازار
روز-پیش، غلبه‌ی شبکه‌های عصبیِ بازگشتی را در بازار درون‌روز، و کارآمدیِ نسبیِ
مدل‌های آماریِ ساده‌تری مانند \lr{LEAR} را در بازار تعادلی گزارش می‌کند
\cite{oconnor_review_2025}.

در کنارِ این یافته‌ها، مرورها به یک کاستیِ روش‌شناختیِ مشترک نیز اشاره می‌کنند که
برای پژوهشِ حاضر اهمیتی محوری دارد. علی‌رغمِ انبوهِ روش‌های پیشنهادی، ارزیابیِ
نظام‌مند و گسترده هنوز کمیاب است: بسیاری از الگوریتم‌ها بر مجموعه‌داده‌های خصوصی و
غیرقابل‌دسترس، در بازه‌های آزمونِ کوتاه و محدود به یک بازار ارزیابی شده‌اند،
به‌ندرت با مدل‌های ساده‌ی مرجع مقایسه شده‌اند، و معناداریِ آماریِ تفاوت‌ها کمتر
آزموده شده است \cite{lago_forecasting_2021, chai_forecasting_2024}. مطالعاتِ
کتاب‌سنجی نیز، در کنارِ ترسیمِ نقشه‌ی انتشاراتِ این حوزه و تمایزِ میانِ
پرارجاع‌ترین و دقیق‌ترین مدل‌ها، بر ضرورتِ اثباتِ دقت از راهِ آزمونِ معناداریِ
آماری تأکید می‌کنند \cite{zema_models_2022}. سرانجام، مرورها کیفیتِ داده،
تفسیرپذیریِ مدل‌ها و تعمیم‌پذیریِ آن‌ها به شرایطِ تازه را از جمله چالش‌های همچنان
گشوده‌ی این حوزه برمی‌شمارند \cite{oconnor_review_2025}. همین کاستی‌ها \lr{—}
نبودِ ارزیابیِ تکرارپذیر و مبتنی بر بنچمارک، و کم‌توجهی به تفسیرپذیری و
تعمیم‌پذیری \lr{—} چارچوبِ انگیزشیِ این پژوهش را شکل می‌دهند و در بخش‌های بعدیِ
همین فصل، به تفکیکِ هر خانواده از روش‌ها، با جزئیاتِ بیشتری بررسی خواهند شد.
